\documentclass[11pt,a4paper]{article}
\usepackage[utf8]{inputenc}
\usepackage[english]{babel}
\usepackage{amsmath}
\usepackage{amsfonts}
\usepackage{booktabs}
\usepackage{hyperref}
\usepackage{listings}
\usepackage{xcolor}
\usepackage{geometry}

\geometry{margin=1in}

\definecolor{codegreen}{rgb}{0,0.6,0}
\definecolor{codegray}{rgb}{0.5,0.5,0.5}
\definecolor{backcolour}{rgb}{0.95,0.95,0.92}

\lstset{
    backgroundcolor=\color{backcolour},
    commentstyle=\color{codegreen},
    keywordstyle=\color{magenta},
    numberstyle=\tiny\color{codegray},
    stringstyle=\color{purple},
    basicstyle=\ttfamily\small,
    breakatwhitespace=false,
    breaklines=true,
    captionpos=b,
    keepspaces=true,
    numbers=left,
    numbersep=5pt,
    showspaces=false,
    showstringspaces=false,
    showtabs=false,
    tabsize=2,
    frame=single
}

\title{Parameter Alignment \& Conflict Report \\ \large Generators vs. \texttt{Clouds.hs} Shader Engine}
\author{Systems Interoperability Group}
\date{\today}

\begin{document}

\maketitle

\section{Executive Conflict Analysis}
A cross-examination of the texture generation scripts (\texttt{generate\_cloud\_noise.py}, \texttt{generate\_blue\_noise.py}) against the dynamic Haskell-Vulkan shader pipeline reveals severe pipeline friction. The structural characteristics (resolution bounds, layout properties, and normalization spaces) baked into the textures directly mismatch the algebraic usage inside the shader core.

---

\section{Detailed Script-to-Shader Discrepancies}

\subsection{1. Channel Pack Symmetry and Frequency Blending}
\begin{itemize}
    \item \textbf{The Generator Output:} \texttt{generate\_cloud\_noise.py} writes out an RGBA 3D texture using independent octaves of Perlin noise in R, and distinct Worley frequencies in G, B, and A.
    \item \textbf{The Shader Call:} Inside \texttt{Clouds.hs}, the erosion math blends channels linearly:
    \begin{equation}
        \text{Erosion} = G \cdot 0.3 + B \cdot 0.15 + A \cdot 0.075
    \end{equation}
    \item \textbf{The Friction:} The Python generator scales frequencies linearly across loops ($f, 2f, 4f, 8f$). However, because the shader applies static weights ($0.3, 0.15, 0.075$), a frequency mismatch arises. If the texture frequency step is too steep, the fine edge details are skipped entirely, manifesting as localized pixel grouping or severe aliasing at structural boundaries.
\end{itemize}

\subsection{2. Blue Noise Dimensionality Mapping}
\begin{itemize}
    \item \textbf{The Generator Output:} \texttt{generate\_blue\_noise.py} generates a flat, $512 \times 512$ 2D spatial texture grid mapping high-frequency distributions onto the $XY$ plane.
    \item \textbf{The Shader Call:} The shader attempts to read this blue noise value to sample structural depth via:
    \begin{lstlisting}[language=Haskell]
    noiseSample <- sample @"blue_noise" uvCoords
    \end{lstlisting}
    \item \textbf{The Friction:} Because the shader wraps the sampler using standard texture filtering over continuous UV projections, the 2D Blue Noise spectrum breaks down due to bilinear reconstruction. This causes the noise properties to collapse into low-frequency blotches when sampled over long-distance horizon rays, preventing effective stochastic layer scattering.
\end{itemize}

\subsection{3. Resolution Sizing Conflict (VRAM vs. Mathematical Density)}
\begin{itemize}
    \item \textbf{The Problem:} The 3D script defaults to a resolution of $128 \times 128 \times 128$. 
    \item \textbf{The Friction:} Combined with \texttt{noiseScale = 0.0008} inside the shader, the texture space repeats its bounds every $1250\text{meters}$ ($1 / 0.0008$). In a raymarching field traversing up to $20,000\text{meters}$ toward the horizon, the shader hits the boundary edges of the $128^3$ texture over 16 times along a single ray, creating repeating cellular rows across the sky.
\end{itemize}

---

\section{Remediation Code Fixes}

\subsection{Fix A: Aligning Python Frequencies with Shader Erosion Scales}
Modify the noise generation setup in \texttt{generate\_cloud\_noise.py} to ensure that octave steps match the geometric scale steps evaluated by the shader erosion passes.

\begin{lstlisting}[language=Python, caption={Corrected Python Frequency Scaling Blocks}]
# Adjust frequencies to pair with shader erosion steps (0.3, 0.15, 0.075)
base_frequency = 4.0
channel_frequencies = {
    'R': base_frequency,        # Base Perlin-Worley Mass
    'G': base_frequency * 2.0,  # Macro erosion structure
    'B': base_frequency * 4.0,  # Medium flaking noise
    'A': base_frequency * 8.0   # High-frequency edge mist
}
\end{lstlisting}

\subsection{Fix B: Shader Texturing Array Adjustments}
To ensure the blue noise pattern functions correctly as a screen-space jitter array rather than blending across geometric space, force integer texel lookups inside \texttt{Clouds.hs}:

\begin{lstlisting}[language=Haskell, caption={Screen-Space Texel Lookup Integration for Blue Noise}]
-- Force direct integer texel extraction to preserve high-frequency jitter fields
let fragCoordX = castFloatToInt (uvCoords.&X * frameWidth)
let fragCoordY = castFloatToInt (uvCoords.&Y * frameHeight)
let texelIndexX = fragCoordX `mod` 512
let texelIndexY = fragCoordY `mod` 512

noiseSample <- readTexel @"blue_noise" (Vec2 texelIndexX texelIndexY)
\end{lstlisting}

---

\section{Pipeline Compatibility Matrix Summary}

\begin{table}[h!]
\centering
\small
\begin{tabular}{llll}
\toprule
\textbf{Parameter Source} & \textbf{Generator State} & \textbf{Shader Expectation} & \textbf{Compatibility Status} \\ \midrule
\textbf{3D Tiling Boundaries} & Strict Linear Wrap & Scaled \texttt{fract()} Lookup & \color{red}Conflict (Triggers repetition) \\
\textbf{Blue Noise Interleave} & 2D Uniform Array & Continuous UV Filtering & \color{red}Conflict (Blurs noise out) \\
\textbf{Erosion Frequency} & Linear $2^n$ Cascade & Hardcoded Scale Weights & \color{orange}Partial Alignment (Needs tuning) \\ \bottomrule
\end{tabular}
\caption{System interface mapping between scripting files and Vulkan pipeline stages.}
\end{table}

\end{document}

